\chapter{Methods}

Before the chapter works through each stage in detail, Figure~\ref{fig:pipeline_overview} summarises the pipeline as a whole: three source corpora are reduced to verbatim expressions, embedded, pooled together with three reference vocabularies into a single shared space, and compared geometrically.

\begin{figure}[h]
\centering
\resizebox{\textwidth}{!}{%
\begin{tikzpicture}[
	node distance=8mm and 12mm,
	box/.style={draw, rounded corners, align=center, minimum height=10mm, minimum width=26mm, font=\small, fill=black!3},
	arr/.style={-{Latex[length=2.5mm]}, thick},
]
	% three source corpora, stacked on the left
	\node[box] (lit) {Scholarly\\literature};
	\node[box, below=3mm of lit] (miv) {MIVILUDES};
	\node[box, below=3mm of miv] (int) {Interviews};

	% main pipeline, left to right
	\node[box, right=of miv] (extract) {Extraction\\(verbatim expressions)};
	\node[box, right=of extract] (embed) {Embedding\\(\texttt{bge-m3})};
	\node[box, right=of embed] (space) {Shared\\cross-corpus space};
	\node[box, right=of space] (analysis) {Geometric\\analysis};

	% reference vocabularies, feeding in from below
	\node[box, below=10mm of embed] (ref) {Reference\\vocabularies};

	\draw[arr] (lit) -- (extract);
	\draw[arr] (miv) -- (extract);
	\draw[arr] (int) -- (extract);
	\draw[arr] (extract) -- (embed);
	\draw[arr] (embed) -- (space);
	\draw[arr] (space) -- (analysis);
	\draw[arr] (ref) -- (space);
\end{tikzpicture}%
}
\caption{Overview of the methodological pipeline. The three source corpora (\ref{sec:vectorising_scholarly_work}, \ref{sec:vectorising_miviludes}) are each reduced to verbatim expressions, embedded, and pooled together with three reference vocabularies (\ref{sec:concept_backbone}, \ref{sec:structural_concepts}, \ref{sec:conceptnet_concepts}) into a single shared space (\ref{sec:shared_space}), which the Results chapter then analyses geometrically}
\label{fig:pipeline_overview}
\end{figure}

\section{Methodological Precedents}


- uncover the semantic structures of what can be concieved as cults
- Vectorizing our understanding cults in the latent space
- in order to do so, I gather different kind of 
- send them into an embeddings model
- perform geometrical comparisons

This matters because it makes your epistemic claim appropriately modest: your embedding configurations are not transparent “maps of cult meaning itself.” They are model-mediated traces of regularities in particular textual or survey-derived inputs

I treat a model’s embedding space as a semantic representation. I do not assume that a vector coordinate corresponds to an independently identifiable semantic feature. Rather, I test whether relational geometry—local neighborhoods, pairwise similarities, category separation, cluster structure, and, where appropriately validated, displacement along constructed directions—captures systematic differences among competing articulations of “cult.”



\footnote{
		The english translation of 'Secte' french substantive IS 'Cult', so when the Mivildues refers to sectarian drifts, it is really talking about cult-like negative consequences. See \ref{sec:defining_cult}
	}



\subsection{Expectations and Biaises}
- I expect to see that that what people consider as cults seems to have more to do with religion than what scholars make it
- I expect to see that the MIVILUDES sectarian drift concept transcends the popular understanding of what a "sect" / "cult" is

\section{Data Collection}

Overall, the goal is to find ways in which the concept of cult can be understood as a collective knowledge. This step constitute the first step of the research project, and is the most important one.
We're looking for associations, definitions, examples, oppositions, limits. In order to achieve this, different type of epsitemologies are gathered.
\subsection{A collection of scholar work about cults}
\subsubsection{Corpus construction protocol}\label{sec:corpus_construction_protocol}
This section describes a reproducible protocol for assembling a corpus of influential books and articles on cults and new religious movements (NRMs). The procedure is designed to be transparent, bias-aware, and replicable: another researcher could rerun the same queries, filters, and ranking rules to obtain a similar set of works.

\paragraph{Data source.} The primary data source is the OpenAlex \autocite{priemOpenAlexFullyopenIndex2022} bibliographic database, accessed via its public API. OpenAlex provides structured metadata for scholarly works, including document type, publication year, and citation counts to measure influence (\verb+cited_by_count+), as well as a topic classification used below for disciplinary stratification.

\paragraph{Search strategy.} To avoid over-reliance on any single label, the corpus is built from a fixed set of keyword queries that combine direct label terms (\enquote{cult}, \enquote{sect}, \enquote{new religious movements}, see \ref{ch:theory}) with the wider \enquote{extreme beliefs} framing; a standalone query for \enquote{sectarian drifts} is added to bridge into the MIVILUDES epistemology (\ref{sec:defining_cult}). The full, fixed query list is reproduced in \ref{sec:keyword_queries} (Table~\ref{tab:keyword_queries_corpus}); each query supports exact phrases (quoted) and boolean AND/OR/NOT/parenthetical combinations. Because OpenAlex stems even quoted single words, \enquote{sect} also matches \enquote{section}; the \enquote{sect} query therefore adds a \texttt{NOT section} clause to suppress that specific false-positive class. \enquote{NRM} is deliberately not searched as a bare acronym alongside \enquote{new religious movements}: it collides with unrelated senses such as \enquote{Natural Resource Management} and \enquote{National Resistance Movement}, both of which fall under the same allowed fields as genuine hits, so no amount of field/subfield screening (below) can disambiguate it — only the unambiguous full phrase is searched. Direct label-term and framing-term queries are matched against the work's \textbf{title} only, rather than title-and-abstract: an abstract-level match only requires the term to appear anywhere in the abstract, which is satisfied by thousands of works that merely mention the term in passing rather than treating it as their subject (e.g. \enquote{cult} alone matches over 12{,}000 works at the title-and-abstract level restricted to the social sciences, dropping to roughly 1{,}000 when restricted to the title, with markedly higher on-topic precision; \enquote{extreme beliefs} similarly drops from 50 hits, mostly generic voting-behavior or wellbeing-psychology papers, to 43 uniformly on-topic hits). Only \enquote{sectarian drifts} remains at the title-and-abstract level, since this framing term is less likely to appear verbatim in a title.

An earlier version of this search paired each direct label term with a disciplinary anchor word in the title query itself (e.g. \enquote{cult} AND sociology, repeated for anthropology, history, psychology, and religion). This was dropped: requiring the literal discipline word in the title is a poor proxy for disciplinary relevance — many on-topic works never use it — and it is redundant with the field/subfield screening described below, which already relies on OpenAlex's own topic classification rather than on title wording. It also produced queries with zero results (e.g. \enquote{sect} AND anthropology, since no anthropology paper happens to have that literal word in its title), needlessly narrowing coverage of otherwise relevant disciplines.

\subsubsection{Screening criteria}
All retrieved records are merged into a single dataset and deduplicated in two passes: first by OpenAlex ID or DOI where available, then by normalized title and first-author surname, which catches near-duplicates that do not share a DOI — reprints, translations, or the same work indexed under several DOIs (e.g. a preprint re-deposited under a new DOI on each revision). Each unique record is then screened against the following predefined inclusion criteria:

\textbf{Type inclusions}: scholarly books, book chapters, and journal articles whose primary subject is cults, sects, or new religious movements, or which offer a central theoretical framing relevant to these phenomena.

\textbf{Type exclusions}: journalistic pieces, activist or apologetic texts, and works where the topic appears only incidentally, unless they are heavily cited theoretical anchors.

\textbf{Field exclusions}: works are also dropped when OpenAlex's own topic classification places them outside the fields relevant to this corpus — retained fields are Social Sciences, Arts and Humanities, and Psychology; works classified under e.g. Medicine, Engineering, or Agricultural and Biological Sciences are excluded automatically, since keyword matches within those fields are reliably incidental rather than substantive.

\textbf{Subfield exclusions}: within the retained fields, the History, Archeology, Classics, and Literature and Literary Theory subfields are further excluded, since \enquote{cult}/\enquote{sect} hits there are reliably about ancient ritual practice or artifacts (e.g. \enquote{the cult of Mithras}, votive cult objects) or about fictional/mythological texts (e.g. Isis and Osiris in Plutarch, Jane Austen's novels) rather than cults as a social phenomenon; Classics is by definition the study of ancient Greek and Roman antiquity, and Literature and Literary Theory by definition analyses fictional or mythological material, so the same reasoning applies to both.

\textbf{Topic exclusions}: works are also dropped when OpenAlex's most granular topic classification — one level finer than subfield — is \enquote{Biblical Studies and Interpretation} or \enquote{Folklore, Mythology, and Literature Studies}, for the same ancient/legendary-content reasoning as the subfield exclusions above; these topics can occur under subfields not otherwise excluded (e.g. Religious studies). More generally, any topic whose name contains \enquote{history} or \enquote{antiquity} is excluded, which generalizes the History subfield exclusion to catch cases where the topic itself is historically framed but the surrounding subfield is not — e.g. \enquote{Mao Cult} is topic-classified as \enquote{Chinese history and philosophy} under the Sociology and Political Science subfield, \enquote{Outcasts, Emperorship, and Dragon Cults in \textit{The Tale of the Heike}} as \enquote{Japanese History and Culture} under Cultural Studies, and \enquote{The Rise of the Hero Cult and the New Simonides} as \enquote{Classical Antiquity Studies} under Anthropology — a fixed subfield blacklist would not catch any of these.

\textbf{Title-pattern exclusions}: works whose title matches \enquote{cult of \ldots}, \enquote{personality cult}, or a fandom/pop-culture or marketing variant (\enquote{cult following}, \enquote{cult classic}, \enquote{cult status}, \enquote{cult collectors}, \enquote{cult phenomenon}, \enquote{cult television}/\enquote{tv}/\enquote{film}/\enquote{movie}/\enquote{cinema}, \enquote{fitness cult}, \enquote{culting}) are excluded, since these all signal the fame/devotion metaphor sense of \enquote{cult} (\cultfame{} in \ref{sec:defining_cult}) — political leader-worship, popular-culture fandom, or brand devotion — rather than the religious/sectarian-group sense this corpus targets (e.g. \enquote{Biography of a Chairman Mao Badge: The Creation and Mass Consumption of a Personality Cult}, \enquote{Cult collectors: nostalgia, fandom and collecting popular culture}, \enquote{The sex lives of cult television characters}, \enquote{...international cult phenomenon: The global expansion of SKAM...}, \enquote{The Culting of Brands: When Customers Become True Believers}). Two further patterns, \enquote{cult(ure)} and \enquote{cult and culture}/\enquote{cults and cultures}, catch titles that typographically or lexically pair \enquote{cult} with \enquote{culture} as a pun (e.g. \enquote{The Confidence Cult(ure)}, \enquote{Between Cult and Culture: Bamiyan, Islamic Iconoclasm, and the Museum}) rather than using the word in its target sense. A final pattern, \enquote{academic sect}, excludes \enquote{sect} used metaphorically for an academic or ideological faction rather than a religious group (e.g. \enquote{Academic Sects as Impediments to Understanding and Progress}).

\textbf{Taxonomic-abbreviation exclusion}: titles containing the pattern \enquote{sect.} (a bare abbreviation followed by a capitalized taxon name) are excluded, since this is the standard botanical and zoological abbreviation for the rank \emph{section} (e.g. \enquote{Sophora sect. Edwardsia}) — a false positive the \texttt{NOT section} clause above cannot catch, since \enquote{section} is never spelled out.

\textbf{Acronym exclusion}: titles containing \enquote{SECT} or \enquote{SECTS} in all capitals are excluded, since an all-caps form reliably signals an acronym rather than the English word (e.g. the \enquote{Scale for Effective Communication in Team Sports (SECTS)}).

\textbf{Ancient-civilization exclusion}: titles naming an ancient civilization, historical period, or the specific ancient religious practice of \enquote{hero cult} (e.g. \enquote{ancient}, \enquote{Roman}, \enquote{Greek}, \enquote{Hellenistic}, \enquote{Homeric}, \enquote{hero cult}, \enquote{Mesopotamian}, \enquote{Byzantine}, \enquote{medieval}, \enquote{Bronze Age}, \enquote{Iron Age}, \enquote{Maya}, \enquote{Aztec}, \enquote{Egyptian}) are excluded. This targets the same ancient-ritual sense of \enquote{cult}/\enquote{sect} as the subfield exclusion above, but at the title level: in practice this sense keeps resurfacing under a different, unpredictable subfield each time the search is rerun (Demography, Communication, Anthropology, and Religious studies have each produced an ancient-cult hit that a fixed subfield blacklist would not catch, e.g. \enquote{Patron Gods and Patron Lords: The Semiotics of Classic Maya Community Cults}), so addressing it directly at the title level is more robust than excluding subfields one at a time. The trade-off is an acknowledged risk of false negatives — a rare modern paper that names an ancient civilization in its title (e.g. a comparative study) could be wrongly excluded — and of false negatives running the other way: the marker list is English-only, so a non-English title can still slip through (e.g. \enquote{Cultes et sanctuaires de l'île de Cos}, caught only on manual review and removed from the approved corpus by hand). Both classes of error require manual review to catch.

\textbf{Publication dates}: research on cults and new religious movements has shifted substantially since the late 20th century, with new theoretical frameworks emerging (e.g. NRM studies displacing earlier \enquote{brainwashing} paradigms, see \ref{ch:theory}). Only works published within the 25 years preceding the final search date are included in the primary corpus: for a search completed in 2026, this corresponds to works published from 2001 onward. Limiting to the last 25 years ensures that the corpus reflects current conceptual debates rather than potentially outdated paradigms. In the API query, this is implemented by filtering on the \texttt{publication\_year} field, e.g. \verb+filter=publication_year:2001-2026+ for a search completed in 2026.

Screening decisions are logged, and excluded records are retained in a separate file for transparency: \texttt{corpus/literature/openalex\_search\_results.txt}, generated by the script reproduced in \ref{sec:openalex_script}. This file also lists the final stratified candidates before manual approval (\ref{sec:corpus_construction_protocol}) folds any of them into the approved literature corpus.

\subsubsection{Ranking and stratification}
After screening, works are separated into two categories: books (including book chapters) and journal articles. Within each category, works are ranked using a mixed influence score that combines:

\begin{itemize}
	\item \textbf{Citation count}: the \verb+cited_by_count+ field from OpenAlex, normalized within each category.
	\item \textbf{Cross-channel presence}: an additional weight is assigned to works retrieved by multiple distinct keyword queries, on the assumption that such works are more central to the field.
\end{itemize}

To ensure disciplinary diversity, the final selection is stratified by disciplinary framing where metadata allows, using OpenAlex's finer-grained topic subfield (e.g. Sociology and Political Science, Anthropology, History, Archeology, Philosophy) rather than its coarser top-level field classification, which is too broad to distinguish these disciplines from one another. For each discipline, the top books and top articles are selected, rather than selecting the top works globally. This reduces the risk that a single disciplinary discourse dominates the corpus while preserving reproducibility.

Within each discipline, only works with some corroborating evidence are eligible for this top selection — either at least one citation, or retrieval by more than one query — unless applying that floor would leave the discipline empty, in which case the unfiltered ranking is used for that discipline only. Without this floor, a discipline with few genuine candidates could otherwise be padded with uncited, single-source records merely to fill the quota (e.g. the same uncited self-published work indexed under several DOIs).

\subsubsection{Final corpus}
The final corpus consists of:
\begin{itemize}
	\item A ranked list of influential books on cults and NRMs.
	\item A ranked list of influential articles on cults and NRMs.
	\item An optional subset of \enquote{bridge texts} that are highly cited across multiple disciplinary framings or retrieval channels.
\end{itemize}

The corpus available for full-text, semantic analysis is a further-conditioned subset of this approved bibliographic list: not every approved record could be obtained as a full-text document — paywalled journal articles and out-of-print books in particular remain out of reach — so the corpus actually analyzed skews toward openly accessible or institutionally reachable works. This access constraint is a source of selection bias distinct from the screening criteria above (\ref{sec:corpus_construction_protocol}), and, like the ancient-civilization marker trade-off described there, is documented for transparency rather than corrected for, given the scope of this thesis.

\subsubsection{Full-text extraction}\label{sec:pdf_extraction}
Once a PDF has been obtained for an approved record, it is converted to clean text for downstream analysis by a small pipeline (\texttt{thesis/corpus/thesis\_corpus/}). Text is extracted with Docling \autocite{auerDoclingTechnicalReport2024}, attempting native extraction first; if the resulting text is implausibly short or empty relative to the page count — signalling a scanned, image-only page — the same page is re-extracted with OCR (Tesseract \autocite{smithOverviewTesseractOCR2007}, English/French/German), so the choice of method is made per document rather than assumed in advance. PyMuPDF is used only to open each file and confirm its page count before conversion, catching corrupted or password-protected files early. Page provenance is preserved throughout: each output retains one record per original page, so any downstream unit of analysis can be traced back to a specific page of a specific source document. Every outcome — including files that could not be opened or whose extraction failed outright — is logged to a per-document status and a corpus-level manifest, rather than silently dropped from the batch.

A successful extraction, however, only establishes that text was obtained, not that it is the intended text. A manual verification pass over the manifest is therefore also part of this stage, and is a distinct source of error from the access constraint above: three documents were found to have extracted cleanly despite being the wrong or an incomplete document — a two-page journal book review standing in for the book itself (Hervieu-Léger, \textit{La religion en miettes}), a several-page fragment in place of a large edited volume (Zablocki and Robbins, \textit{Misunderstanding Cults}), and a paywalled journal article of which only the citation/abstract page had been saved (Prohl, \textit{Enjoying Religion}, saved twice under different filenames). All three were removed from the corpus (source files, manifest entries, and bibliography records alike) once identified. No automated check in the pipeline distinguishes a technically-successful extraction of the wrong document from a correct one; this remains a manual, spot-checked judgment.

\subsubsection{Named cults and dimensions}
From the Corpus gathered about, the goal is to extract named entities of cults. Of course these entities can also be mentioned with doubts, there should also be extract the certainty at which the entity is pointed at as a cult by the literature. 
Similarly we want to do the same with named dimensions of cults, with a certainty ammount. 

\subsection{MIVILUDES}

\subsubsection{Corpus}
With the MIVILUDES we have access to their public information on their website.
\Textcite{CommentIdentifierDerive} already provides a definition attempt of a sectarian drift, as well as a list of these drifts; both are extracted here.
\Textcite{QueDitLoi} sets out the implications under French law, and more usefully, a list of legal infractions that sectarian drifts are likely to commit; these are extracted here as well.

Additionally, \textcite{missioninterministerielledevigilanceetdeluttecontrelesderivessectairesmiviludesRapportDactiviteMiviludes2024} provides extensive information about MIVILUDES's own work.


\subsubsection{Named cults}
- 
\subsubsection{Dimensions}



1. Corpus
	1.0 Textual data (not published, just the bibliography)
	1.1 Named entities - cults mentioned and their level of certainty or opposition
	1.2 Cult dimensions - cults mentioned and their level of certainty or opposition

2. MIVILUDES Content and reports.
	1.0 Textual data (not published)
	2.2 Named entities that were trialed from the work of the MIVILUDES
	2.3 Legal basis for conviction
3. A survey
	3.1 Named entities - cults mentioned and their level of certainty or opposition
	3.2 Cult dimensions - cults mentioned and their level of certainty or opposition

4. Other data
 	4.1 Cult databases 
		4.1.1 List from Wikipedia
	4.2 Dctionaries
		4.2.1 Most used terms in cultural studies
		4.2.2 Most used nouns in the english language
	4.3 Other
		4.3.1 Mythemes




2. A collec

First and foremost, we gather a corpus of scholar texts about the subject. 

In order to gather a lot of relevant litterature efficiently, open academic APIs can be queried.

Queries were designed to cover major theoretical dimensions of cults and new religious movements (definition, organization, conversion, authority, boundaries etc.) as identified in different fields.




\subsection{Scholarly definitions}


  - Selected corpus of cult studies / New Religious Movement (NRM) studies
    -- methodology for corpus selection.
  - Selected corpus of sectarian drifts (MIVILUDES).
  - Prototype gathering (survey / free-listing).
- Gather data from dictionaries to test different things:
  - Naive, a priori dimensions.
  - A dictionary of most common substantives in English.
  - Religious studies vocabulary.
  - Other semantic data sets worth comparing against. (mythemes)


\subsection{MIVILUDES Sectarian drifts}
\subsection{Survey}
- First I had to create a methodology for my survey \ref{ch:survey_methodology}. 
- Explain why I chose this methodology and what were the aims.
- I decided to use the same methodology on instagram, by posting in my story (insert screenshot).
- I ran the first methodology \ref{ch:survey_methodology} for gathering the exemplars, open ended questions in real life. 
- The interview gave definition attempts and examplars.
- I used google record \autocite{google_recorder_2026} which uses automatic speech recognition to transcribe, then I used an LLM (Claude 5 Opus) to batch process the transcriptions "/Transcripts edited" for the full transcripts fix transcription errors and extract the demographic data. 
This data is the extracted into a json file with the following:
	- Demographic data (age, gender, nationality, main language, country of residence)
	- Cult exemplars in order of mention
- Then I did a manual pass, listening to the recording and making sure the transcription was accurate and the data was correctly captured, and anonymised.


\section{Data Extraction}

\section{Vectorisation}

\subsection{Overview of Gathered Data}\label{sec:vectorisation_data_overview}
The table below summarises every dataset this stage of the thesis operates on, before the following subsections describe how each is actually processed. Three are corpora of expressions extracted from full text (\ref{sec:vectorising_scholarly_work}); four are small, curated or derived reference/emergent lists (\ref{sec:vectorising_miviludes}, \ref{sec:concept_backbone}, \ref{sec:structural_concepts}, \ref{sec:emergent_entities}) rather than extracted directly from full text; the last is derived from all seven of the others combined (\ref{sec:shared_space}). All embeddings, throughout, are produced by the same local model (\texttt{bge-m3} \autocite{chenM3EmbeddingMultiLingualityMultiFunctionality2025}, 1024-dimensional, multilingual, served via Ollama) — no dataset uses a different embedding model from any other.

\begin{center}
\begin{tabularx}{\textwidth}{@{}lXX@{}}
\toprule
Dataset & Size & Storage \\
\midrule
Scholarly literature corpus & 57 documents, 39{,}236 extracted expressions & \texttt{processed/literature/criterion\_expressions.jsonl} \\
MIVILUDES corpus & 2 documents (the 2024 activity report; the \enquote{Comment identifier} reference text), 914 extracted expressions & \texttt{processed/miviludes/criterion\_expressions.jsonl} \\
Interview corpus & 26 transcripts, 230 extracted expressions & \texttt{processed/interviews/criterion\_expressions.jsonl} \\
MIVILUDES official criteria & 17 criteria, 17 embedded points (French only; English retained as a display label) & \texttt{metadata/miviludes\_criteria\_embedded.jsonl} \\
Generic concept backbone & 3{,}000 WordNet concepts & \texttt{dictionaries/concept\_backbone\_embedded.jsonl} \\
Structural concepts & 1{,}500 corpus-derived generic/relational terms & \texttt{dictionaries/structural\_concepts\_embedded.jsonl} \\
Emergent entities & 3{,}251 named entities/dimensions mentioned $\geq$3 times across all corpora & \texttt{processed/shared\_space/embedding\_space.jsonl} (\texttt{source\_dataset=emergent\_entities}) \\
Shared cross-corpus space & 44{,}325 points pooled from all seven datasets above (after removing duplicate/short expressions, see \ref{sec:shared_space}), reduced to 394 dimensions & \texttt{processed/shared\_space/embedding\_space.jsonl} \\
\bottomrule
\end{tabularx}
\end{center}

Every corpus above shares the same per-expression schema: a verbatim source quote, a short self-contained embedding text, named entity/dimension anchors, and three descriptors of how the expression makes its claim (claim mode, epistemic status, attribution) — detailed in full in \ref{sec:vectorising_scholarly_work}. A separate, machine-generated index (\texttt{processed/registry.csv}) tracks, for every document across all three corpora, whether it has been extracted and embedded yet — not itself analytical data, but the bookkeeping that keeps the other three corpora's processing state auditable as they grow. The figures in the table above, and the shared space described in \ref{sec:shared_space}, are built from the first-generation extraction pipeline (\ref{sec:vectorising_scholarly_work}) and are retained on disk as a frozen developmental baseline, not modified further. A revised extraction design (\ref{sec:extraction_v2}) supersedes it; the shared space rebuilt from the revised archive, once extraction of the full corpus under that design was complete, is described in \ref{sec:shared_space_v2}, with the corresponding revision to emergent-entity construction in \ref{sec:emergent_entities_v2} — every geometric result in this thesis from that point forward is computed against the revised space, not the one summarised in the table above.

\subsection{Vectorising the MIVILUDES Framework}\label{sec:vectorising_miviludes}
Unlike the scholarly corpus, the MIVILUDES framework supplies its own dimensions directly: the Miviludes has published a fixed list of seventeen criteria for identifying a \enquote{dérive sectaire} (\textit{Comment identifier une dérive sectaire?}). Each criterion is embedded with the same local model as every other corpus (\texttt{bge-m3}); although the source text exists in both French and English, only the French version — the official original — is embedded and pooled into the shared space (\ref{sec:shared_space}). The English translation is retained purely as a display label alongside the same point, not as a second geometric point: embedding both would not add analytical value beyond confirming translation fidelity, which is instead checked directly, once, via cosine similarity between the two raw (pre-reduction) embeddings for all seventeen criteria — mean 0.87, minimum 0.50 (\textit{L'importance des démêlés judiciaires} / \enquote{The extent of the group's legal entanglements}), both of the lowest-scoring pairs inspected by hand and found to be accurate translations rather than errors, consistent with a general pattern of short official phrases embedding less closely across languages than their length alone would suggest.

A further, not-yet-built use of this framework remains open for later work: scoring named entities tried under French sectarian-drift law against these same seventeen criteria, e.g. representing entity $A$ as $(1,0,0,\ldots)$ if only the first criterion was found to apply to it, entity $B$ as $(0,1,0,\ldots)$, and so on — a criterion-membership representation distinct from, and not requiring, the embedding-based one described above.

\subsection{Vectorising Scholarly Work}\label{sec:vectorising_scholarly_work}
Unlike the MIVILUDES framework (\ref{sec:vectorising_miviludes}), whose dimensions are given directly by the source, the scholarly corpus offers no single agreed-upon definition of a cult to anchor a vector space to in advance. The approach taken here is instead to first identify, within each approved document's full text (\ref{sec:pdf_extraction}), the individual expressions through which the literature articulates criteria for cult-like or sectarian formations, and only embed those expressions afterward — rather than assuming the relevant dimensions ahead of time.

Each document's page-level text is split into coherent chunks of roughly 300--700 words (or 2--5 paragraphs), keeping page provenance so every chunk remains traceable to a specific page of a specific source. Each chunk is passed to a locally-run large language model (\texttt{qwen3:4b} \autocite{yangQwen3TechnicalReport2025}, served via Ollama) instructed to identify expressions that describe, attribute, dispute, reject, define, or problematise cult/sectarian criteria — explicitly not to judge whether anything described is objectively a cult. This distinction matters for the same reason the screening criteria in \ref{sec:corpus_construction_protocol} exclude ancient and metaphorical senses of \enquote{cult}: the corpus is asked what criteria are \textit{articulated}, by whom, and with what degree of certainty, not what the analysis itself endorses as true. Each extracted expression is accordingly tagged with its claim mode (a direct statement, an attributed or quoted claim, a definition, or a reflective question), its epistemic status (asserted, qualified, contested, negated, or speculative), and its attribution (the author, a cited author, a participant, an institution, or a journalist) — so a claim a text disputes is recorded distinctly from one it asserts.

The same pipeline was also run over MIVILUDES's own non-framework source text (a reference document distinct from the seventeen-criteria list itself, \ref{sec:vectorising_miviludes}) and over the interview transcripts (\ref{ch:survey_methodology}), via two small adapters that convert non-PDF sources — a bare document and a batch of transcripts, respectively — into the same page-provenance structure the pipeline otherwise expects from a PDF. One consequence specific to the interview transcripts: the model would sometimes echo the transcript's own literal speaker labels (\enquote{Interviewee:}, \enquote{Interviewer:}) into the attribution field rather than mapping them onto the categories above; this was corrected by normalizing known synonyms (e.g. \enquote{interviewee} to \enquote{participant}) before validating the model's output, rather than relying on prompt wording alone to constrain a small local model. A related, separate issue affected 3 of the 26 transcripts: a trailing transcriber's-notes section (documenting ASR corrections and uncertain names, not spoken by either party) was being chunked and annotated as though it were dialogue — traced to one extracted item (\enquote{cercle solaire}) that turned out to quote the note itself rather than the interviewee, and fixed by excluding this section before chunking; inspecting the affected items by hand also surfaced and corrected several interviewer-question/participant-answer attribution mislabels elsewhere in the interview archive.

Each extracted item's position within its own interview (chunk order, then extraction order within a chunk) is preserved as \texttt{response\_rank}, carried through unchanged into the shared space (\ref{sec:shared_space}) — but this is extraction-order provenance only, not a free-listing rank or a cognitive-salience proxy. The interview protocol (\ref{sec:interview_protocol}) is not a multi-item free-listing task: one prompt (\enquote{what comes to mind when you hear the word cult?}) elicits a single example, and a second prompt together with open-ended probes ask the participant to justify \textit{that same example}, not to list further items. A direct check (\texttt{thesis\_corpus.audit\_free\_listing\_rank}) confirms \texttt{response\_rank} cannot support a salience claim even at its narrowest, most favorable reading — whether \texttt{response\_rank == 1} is the participant's own first claim rather than, for instance, the interviewer's own question text — which holds for only 11 of 26 interviews. Geometric analysis of the interview side accordingly uses a separate, manually-reviewed initial-exemplar workflow instead: each interview's first participant claim is identified by a proposal heuristic, checked by hand against the original transcript, and only then analyzed geometrically (\texttt{thesis\_corpus.propose\_initial\_exemplars}, \texttt{analyze\_initial\_exemplars}) — never \texttt{response\_rank} itself.

To preserve short but complete initial interview responses that the general expression-quality filter excludes (\ref{sec:shared_space}), a separate, manually-reviewed prototype layer was constructed rather than relaxing that filter itself. Each verified opening response was resolved against the unfiltered interview archive, using its existing embedding, and projected through the fitted scaling and PCA transformation of the already-established shared space (\texttt{thesis\_corpus.build\_interview\_prototype\_layer}). Prototype points therefore occupy the same coordinate system as the pooled data without modifying the fitted space or reintroducing short fragments into the general interview-expression corpus. Of the 26 interviews, 25 yielded a reviewable and projectable opening response; one was retained as unavailable because no qualifying participant claim could be identified in that transcript at all.

Every extracted expression, along with the named entities and dimensions (\enquote{entity anchors}) it mentions, is embedded with a second local model (\texttt{bge-m3}, 1024-dimensional, multilingual). Running both models locally via Ollama rather than a cloud API keeps the corpus's not-yet-published full texts off any third-party service and keeps the pipeline reproducible independent of an external API's availability or pricing.

Because the resulting archive keeps every field — including full chunk context and a vector per entity anchor, for traceability — it is large and grows with the corpus. A reduced working file is therefore derived from it for interactive analysis and visualisation: large or redundant fields are dropped, and the set of expressions is optionally downsampled by projecting embeddings to a lower dimension (PCA) and clustering them (MiniBatchKMeans \autocite{sculleyWebscaleKmeansClustering2010}, via \texttt{scikit-learn} \autocite{JMLR:v12:pedregosa11a}), sampling a fixed number of expressions per cluster rather than truncating arbitrarily — keeping the working file a representative, size-bounded sample of the full archive rather than, for instance, only its first rows.

Early testing surfaced two further extraction characteristics, addressed at pooling time rather than at extraction (\ref{sec:shared_space}) in this first-generation archive: the model occasionally emits the same expression twice within a single chunk's response, and occasionally emits a short, low-content phrase — already captured as an entity anchor on a fuller expression — as though it were an independent expression in its own right. Both recur, in different form, among the problems a later, larger fidelity review found at greater scale (\ref{sec:extraction_v2}), and both are addressed structurally, at extraction time itself, in the revised pipeline described there.

\subsection{Revising the Extraction Pipeline: Verbatim Spans and Automated Verification}\label{sec:extraction_v2}

The pipeline described above (\ref{sec:vectorising_scholarly_work}) asked the extraction model for two related but distinct things per identified expression: a verbatim \texttt{source\_quote} substring of the chunk, checked programmatically against the chunk text, and a separate, model-produced \texttt{embedding\_text} intended to be a short, self-contained restatement of it — the field this thesis actually embeds and analyses. Only the first of these two fields was ever verified. A manual fidelity review of a stratified 100-item sample drawn directly from the three raw extraction archives, conducted after the corpus described above had already been built, found this unchecked second step to be the pipeline's principal weakness: examples included expressions that were off-topic despite containing a cult-related word, expressions spanning several unrelated claims, short conversational fragments with no standalone content (e.g.\ \enquote{Well, perhaps to some extent.}), section headings and cited work titles extracted as though they were assertions, and — the most direct evidence of the underlying mechanism — spans visibly truncated or silently rewritten relative to their own source chunk.

Reconstructing the relationship between \texttt{embedding\_text} and its own source chunk across the full literature archive showed why: only 17\% of embedding texts were identical to their verified \texttt{source\_quote}, 51\% were a shorter span drawn from within it, and the remaining 32\% were not a substring of the source chunk at all. The manual review's clearest failures split cleanly along this line. Truncations that cut a quotation before its meaning was complete — a quoted sentence stopped mid-clause, silently reversing the sense of what it went on to say — came from the sub-span category: a legitimate editorial trim in most cases, but one the pipeline had no way to distinguish from a meaning-altering one. Outright rewrites, mistranslations, and garbled substitutions (e.g.\ a French sentence's date, \enquote{2009}, silently rendered in the English \texttt{embedding\_text} as the clock time \enquote{2:09}) came from the third category — expressions the model had generated rather than copied.

The extraction pipeline was accordingly redesigned around one governing rule: \texttt{embedding\_text} is never produced by the model. Each chunk is still sent to a locally-run model in a single structured call (\texttt{qwen3:4b} \autocite{yangQwen3TechnicalReport2025}, JSON-schema-constrained output via Ollama, one retry on a malformed response before the chunk is logged as failed and contributes nothing, rather than a guessed or partial result), but the model now returns only a \texttt{verbatim\_expression} field, which is resolved programmatically against the same Unicode-normalised chunk text the model was shown; any candidate that does not resolve to an exact, contiguous substring — a paraphrase, a translation, a truncation — is discarded before it ever reaches the embedding step. The revised prompt is also deliberately conservative: it asks the model to keep at most four expressions per chunk, ranked strongest first, and explicitly to prefer returning nothing over returning something doubtful, reversing the earlier design's implicit incentive to extract as much as possible from each chunk.

A verbatim span passing this check is then subject to a second, fully deterministic layer of screening — ordinary code, not a further model call — before it is accepted. This layer rejects, without exception: text overlapping a region the source PDF is independently known to have corrupted (below), a documented ligature-substitution pattern specific to one affected source, headings and page furniture, bibliographic entries and cited work titles, standalone personal names, citation-dominated spans, and — for the interview corpus — any span attributed to the interviewer. It separately rejects short spans (six words or fewer) unless they name a specific entity or concept, or are explicitly tagged as a self-contained lay association occurring in a passage clearly about cults or sects; this distinction was added deliberately so that short but meaningful lay expressions (e.g.\ a participant's \enquote{Tomato cult!}, or \enquote{la Scientologie}) are not discarded merely for being short — a distinction the first-generation pipeline's plain length-based pooling filter (\ref{sec:shared_space}) could not draw. Milder textual irregularities — a detached accent glyph, an unusual script, an unbalanced quotation mark — are recorded as flags for later inspection rather than treated as grounds for automatic rejection, so that this screening layer does not itself become a second, uninspectable source of silent loss.

\textbf{A text-integrity audit of the source PDFs.} Tracing the manually-reviewed sample's garbled examples (e.g.\ \enquote{di Y cult} for \enquote{difficult}) back through the pipeline showed the corruption was already present in the very first extraction step (\ref{sec:pdf_extraction}), not introduced by the language model: certain PDFs' embedded fonts encode common ligatures (\textit{ffi}, \textit{ffl}) as separate glyphs that the extraction step decodes as an isolated capital letter, and three PDFs' character-mapping tables are broken outright, substituting an unrelated Unicode block (Latin Extended-B, Greek, or Private-Use-Area codepoints) for ordinary Latin text over parts of the document — a defect the pipeline's own quality check (\ref{sec:pdf_extraction}), which only counts characters, cannot detect, since a substitution cipher preserves the expected character count exactly. A dedicated, read-only audit script scores every document's extracted text for these and related defects (replacement characters, control characters, detached accent glyphs, mixed-script tokens) and records, per document, which of these apply; the extraction screen above consults this record directly rather than re-deriving it, so a single audit run is the sole source of truth for which documents and passages are treated as textually compromised.

\textbf{Interview attribution, revised.} The first-generation pipeline's interviewer/participant mislabelling (\ref{sec:vectorising_scholarly_work}) was corrected by hand, case by case, but nothing prevented the same class of error from recurring. The revised pipeline instead parses each transcript into speaker turns from the transcript's own labels (\enquote{Interviewer:}, \enquote{Interviewee:}, and a transcriber's-notes marker where present) before any expression is accepted: an extracted span's attribution is set by which turn it falls inside, never by the model's own claim, and a span crossing a turn boundary, containing a speaker label itself, or falling in the transcript's demographic header is rejected outright. This closes the gap structurally rather than by further manual correction.

\textbf{Automated verification in place of full manual review.} Reviewing every retained expression of a corpus this size by hand was not feasible within this thesis's timeline, so a second, independent verification step was added directly into the pipeline rather than left to the researcher's own later spot-checking alone. Every expression the deterministic screen retains is shown, together with its full source chunk and its assigned labels, to a second, separately-prompted local model — a larger instance of the same model family used for extraction (\texttt{qwen3:8b} \autocite{yangQwen3TechnicalReport2025}) — which answers the same faithfulness, self-containedness, cult-relevance, textual-intelligibility, and atomicity questions a human reviewer would be asked, plus whether each assigned label (attribution, claim mode, epistemic status) looks correct and whether a clearly better formulation of the same idea exists elsewhere in the chunk (flagged for review rather than substituted automatically, since inserting a different span in its place would itself violate the verbatim-extraction rule above). A fixed rule, not a further judgement call, accepts an expression only if every one of the five quality checks is affirmed and no problem is named; a verification call that fails outright rejects the expression rather than letting it through unverified. Every verdict — accepted or not — is retained in full alongside the expression it concerns, and is treated throughout this thesis as \textit{model-judged}: a documented, automated signal, never a substitute for human review or evidence about the extracted content's substantive accuracy. Before being applied to the full corpus, this design was piloted on a small, deliberately mixed sample of chunks — spanning cleanly-extracted and known problematic source documents alike, including several of the specific passages that had produced the garbled or over-broad expressions described above — and an independent, code-level check confirmed that none of those specific known failures reappeared among the expressions the revised pipeline and verification step together retained.

The first-generation archive described in \ref{sec:vectorising_scholarly_work}, together with the shared space built from it (\ref{sec:shared_space}) and the corrections documented in \ref{sec:known_limitations}, is retained on disk as a frozen developmental baseline and is not modified further; the corpus this thesis analyses from this point forward is produced by the revised pipeline described here.

\subsection{A Generic Concept Backbone}\label{sec:concept_backbone}
Each corpus-specific vectorisation above (\ref{sec:vectorising_miviludes}, \ref{sec:vectorising_scholarly_work}) anchors its embedding space to dimensions drawn from that corpus itself — the MIVILUDES framework's own criteria, or expressions extracted from the scholarly literature. Comparing spaces built this way risks comparing artefacts of each corpus's own vocabulary rather than a shared underlying geometry. To give every corpus a common, corpus-independent reference frame, a separate set of anchor points was built: a list of general, abstract English concepts, selected without reference to cults, sects, or any term specific to this research topic, so that its influence on the resulting space is structural rather than thematic.

The source is the Open English WordNet \autocite{mccraeEnglishWordNet20192019} (\texttt{oewn:2024}), accessed via the \texttt{wn} Python library \autocite{goodmanIntrinsicallyInterlingualWn2021}.\footnote{ConceptNet \autocite{speerConceptNet55Open2018} was considered first for this topic-neutral backbone specifically, being multilingual and relation-based rather than a flat word list, but was set aside for that particular purpose: its public API has been persistently unavailable, and its bulk data dump does not encode a shared identifier across languages — direct English-to-French edges for common abstract words were absent in a partial scan of the dump, which would have made a multilingual anchor list unreliable without ad-hoc heuristics. ConceptNet's English-only bulk dump was used later for a different, deliberately non-neutral purpose — generalising the corpus-derived structural-concepts vocabulary below (\ref{sec:conceptnet_concepts}) — where those same limitations don't apply.} The construction is purely structural, with no hand-picked seed terms at any step. Starting from every noun synset in the lexicon (84,956 in total), a synset is kept only if its hypernym chain leads to the WordNet root synset \textit{abstraction, abstract entity} rather than to \textit{physical entity} or another root — WordNet's own top-level split between abstract and concrete nouns (41,351 pass). Synsets denoting named entities are then excluded, identified either by an \texttt{instance\_hyponym} relation (WordNet's mechanism for individuals, e.g. \enquote{Paris}, as opposed to ordinary class hyponymy such as \enquote{city}) or by a capitalised primary lemma (9,151 excluded). A further exclusion removes biological taxonomy terms (lexicographer files \texttt{noun.animal} and \texttt{noun.plant}, e.g. \enquote{bird genus}, \enquote{asterid dicot genus}): these pass the abstractness test structurally, since a taxonomic genus is a hyponym of \textit{group}, but their hyponym counts are inflated by simply enumerating species, which dominated early rankings with domain-specific taxonomy rather than general concepts (875 excluded). Each of the 31,325 remaining candidate synsets is scored by node degree — the total number of WordNet relations (hypernym, hyponym, meronym, holonym, attribute, \textit{similar\_to}, and others) touching it — and the top-ranked 3,000 are kept as the concept backbone, each retained with its WordNet gloss for sense disambiguation before embedding (a bare polysemous word such as \enquote{state} is a poor anchor without the sense it names). This yields general concepts such as \enquote{law}, \enquote{chemistry}, \enquote{biology}, \enquote{mathematics}, \enquote{quality}, and \enquote{time period} rather than corpus-specific vocabulary. The resulting list is embedded with the same model as the rest of the corpora (\texttt{bge-m3}) and used as a fixed set of anchor points when comparing spaces built from different corpora.

\subsection{Structural Concepts: a Corpus-Derived Reference Subset}\label{sec:structural_concepts}
The concept backbone's topic-neutrality is exactly what makes it useful as an independent yardstick, but it comes at a geometric cost: measured in the shared space (\ref{sec:shared_space}), its centroid sits notably farther from the corpus-expression centroid than emergent entities' does (16.7 vs.\ 12.5 shared-space units, against a roughly 33-unit expression-to-expression baseline) — a real, if moderate, effect, plausibly a register mismatch between formal WordNet glosses and corpus prose. Rather than replacing the concept backbone, a second \texttt{point\_role="reference"} subset was added alongside it, built the opposite way on purpose: extracted \textit{from} the corpora's own expression text, so it sits closer to the data by construction, at the cost of no longer being topic-neutral — it exists because the corpora use this vocabulary. The two subsets serve different questions: the concept backbone when independence from the corpus matters, structural concepts when interpretive proximity matters more.

The obvious source — pooling from the entity anchors already used for emergent entities (\ref{sec:emergent_entities}) — was tried first and abandoned: entity anchors capture concrete named things (\enquote{Scientology}, \enquote{charismatic leader}), not abstract social-structure vocabulary. Summed across the entire corpus with no threshold at all, words such as \enquote{control} (4 mentions), \enquote{authority} (12), \enquote{manipulation} (3), and \enquote{harm} (0) barely register as tagged anchors, at any frequency cutoff. Instead, the actual expression prose is tokenized directly, and a token is kept only if it is a valid Open English WordNet lemma (the same lexicon behind \ref{sec:concept_backbone}, so no new resource) falling in an in-domain lexicographer file — nouns and adjectives from social, relational, cognitive, and group-dynamics domains, excluding concrete/physical ones (places, objects, body parts, dates, quantities) that pass plain WordNet membership but aren't structural concepts in the relevant sense; verbs are excluded entirely, since the target vocabulary names a concept or role rather than an action. WordNet membership itself does double duty here: it discards proper nouns, typos, and non-English tokens (MIVILUDES is mostly French) essentially for free, and supplies a ready-made gloss for every survivor, per the same \enquote{\textit{term}: \textit{gloss}} embedding convention as the concept backbone. A word is counted once per expression containing it, not once per raw occurrence, so a single repetitive sentence cannot inflate its rank. A small number of high-frequency words whose top-ranked WordNet sense was clearly wrong for how the corpora use them (e.g.\ \enquote{religious} defaulting to a noun sense meaning a monk or nun, rather than the adjective) were hand-corrected, and a handful of remaining proper-noun leaks and mis-sensed matches were spot-checked and excluded. This surfaced a systematic gap rather than a one-off: Open English WordNet's \texttt{instance\_hyponym} marking turns out to be incomplete for minor historical figures, so common words coinciding with their entries (e.g.\ \enquote{smith}, \enquote{land}, \enquote{king}) leaked through, including two on-topic-but-wrong leaks that would have undermined the dataset's own purpose — \enquote{hubbard} (L.\ Ron Hubbard) and \enquote{iskcon} (a specific named sect, the Hare Krishnas) are exactly the kind of named-person/named-group contamination this filter exists to keep out, structural concepts being meant to describe roles and dynamics rather than name specific people or groups (that job belongs to emergent entities, \ref{sec:emergent_entities}). A systematic sweep for biographical-looking glosses (matching patterns such as \enquote{United States \ldots}, a birth year, or a year range) caught the bulk of these at once; verified zero such matches remain in the final list. Neither correction is exhaustive — a 1{,}500-entry automatically ranked list carries residual noise beyond what a spot check catches, documented here rather than claimed fully clean, in the same spirit as the two known-but-deferred extraction issues already noted in \ref{sec:vectorising_scholarly_work}.

Of 33{,}113 unique post-stopword tokens found in the expression text, 7{,}053 survive WordNet and domain filtering; the top 1{,}500 by expression-frequency are kept as structural concepts. This size was chosen empirically rather than arbitrarily: mentions per term fall off gradually (70 at rank 600, 40 at rank 1{,}000, 24 at rank 1{,}500), and proper-noun leakage and per-term mention counts both degrade noticeably past that point, so 1{,}500 sits at the practical ceiling before quality drops rather than being picked in advance. The most frequent are genuinely structural rather than named: \enquote{religious} (2{,}688 mentions), \enquote{movement} (1{,}512), \enquote{church} (1{,}250), \enquote{cult} (1{,}229), \enquote{group} (1{,}129), and \enquote{religion} (1{,}066) — unlike emergent entities' own top mentions, which are named groups.

\subsection{A ConceptNet-Generalised Third Reference Vocabulary}\label{sec:conceptnet_concepts}
The concept backbone (\ref{sec:concept_backbone}) and structural concepts (\ref{sec:structural_concepts}) sit at two ends of a trade-off: fully topic-neutral versus fully corpus-derived. A third \texttt{point\_role="reference"} subset sits between them — generalising the structural-concepts vocabulary via ConceptNet's \autocite{speerConceptNet55Open2018} associative graph, rather than drawing new terms directly from corpus text or from a topic-neutral yardstick. Seeded from the 1,500 structural-concept terms, every English ConceptNet edge touching one of them was retrieved from the bulk assertions dump (475,700 matching edges), excluding loose or morphological relation types (\texttt{RelatedTo}, which alone accounts for 63\% of all matched edges at the lowest average edge weight of any major relation type, plus \texttt{HasContext}/\texttt{DerivedFrom}/\texttt{FormOf}/\texttt{EtymologicallyRelatedTo}) and oppositional ones (\texttt{Antonym}/\texttt{DistinctFrom}/\texttt{NotDesires}). Candidate terms were ranked by \textit{specificity} — the fraction of a term's total ConceptNet connectivity accounted for by the structural-concepts seeds, rather than raw seed-edge weight — since ranking by raw weight alone floods the result with generic hub words (\enquote{make}, \enquote{have}, \enquote{activity}) connected to nearly everything in ConceptNet, seed or not. The same in-domain Open English WordNet lexicographer-file filter used for structural concepts (\ref{sec:structural_concepts}) was applied again, on the same reused code path.

This automated pipeline still surfaced residual hub-word noise near the top of its ranked candidates, so every one of the 1,345 resulting candidates was manually reviewed — the same hand-curation step already applied to structural concepts' own automatically-ranked list. 195 were kept as genuinely relevant to the domain (e.g.\ \enquote{spell}, \enquote{oracle}, \enquote{mentor}, \enquote{grooming}, \enquote{torah}, \enquote{yoke}); 1,150 were flagged as generic and excluded, never embedded or pooled. Only the 195 kept terms are embedded (same \enquote{\textit{term}: \textit{gloss}} convention as the other two reference vocabularies) and pooled into the shared space (\ref{sec:shared_space}) as \texttt{source\_dataset="conceptnet\_concepts"}.

\subsection{Emergent Entity Anchors}\label{sec:emergent_entities}
Every extracted expression above (\ref{sec:vectorising_scholarly_work}) is tagged with the named entities and dimensions it mentions (\enquote{entity anchors}, e.g. \enquote{Scientology}, \enquote{charismatic leader}), and each one is individually embedded at extraction time alongside the expression itself. Until now this sat unused beyond the expression it was attached to. To let named entities be compared geometrically — both to the corpus expressions that mention them and to the topic-neutral concept backbone (\ref{sec:concept_backbone}) — these per-anchor embeddings are pooled into a set of points of their own, called \textit{emergent entities}: named entities/groups/concepts that emerge \textit{from} the corpora themselves, as opposed to an expression point (a claim a source makes) or a reference point (the corpus-independent concept backbone) — see \ref{sec:shared_space} for the full three-way distinction.

Anchors are normalised (lowercased, whitespace-collapsed) and pooled once across all three corpora together, so the same entity mentioned in the literature, MIVILUDES, and an interview becomes a single point. Of 19{,}596 unique normalised anchors, most are mentioned only once or twice — mostly noise (typos, overly specific one-off phrases) rather than recurring entities; only anchors mentioned at least three times across the whole corpus are kept, giving 3{,}251 emergent-entity points, deliberately the same order of magnitude as the 3{,}000-entry concept backbone. The most frequently mentioned are \enquote{Scientology} (3{,}867), \enquote{charismatic leader} (3{,}417), \enquote{NRMs} (1{,}279), and \enquote{Heaven's Gate} (577). The full ranking, top 100 with a per-corpus mention breakdown, is reproduced in \ref{ch:emergent_entities}.

\subsection{Revising Emergent-Entity Construction}\label{sec:emergent_entities_v2}

Rebuilding the emergent-entity layer under the revised extraction pipeline (\ref{sec:extraction_v2}) at first produced an implausible result: only 21 entities survived the same construction described above (\ref{sec:emergent_entities}), almost all mentioned exactly once, none shared across all three corpora. Tracing the cause showed this was not a fact about the corpus but a consequence of where entity anchors come from: the extraction model tags anchors only on expressions that survive the full verbatim-span and deterministic-screening pipeline (\ref{sec:extraction_v2}), so once that pipeline began retaining only 14--28\% of what the first-generation pipeline kept per corpus (\ref{sec:vectorisation_data_overview}), the same proportion of named-entity mentions sitting inside now-rejected spans was silently discarded along with them — including mentions inside text that was itself unremarkable, just not judged a standalone claim worth keeping.

The fix draws on a second field the extraction model already produces per chunk but that had never been pooled into the shared space: \texttt{domain\_terms}, explicitly defined at extraction time as \enquote{cult-related concepts or named groups that appear in the text, even if no expression around them is worth keeping} — a deliberately broader net than entity anchors, checked directly against the literature corpus alone: 52{,}179 unique \texttt{domain\_terms} (100{,}031 mentions) against the 20 entity-anchor-only candidates the unrevised construction found there. This breadth comes at a cost: \texttt{domain\_terms} mixes genuine named entities (\enquote{Scientology}, \enquote{Heaven's Gate}) with generic domain vocabulary (\enquote{cults}, \enquote{brainwashing}, \enquote{mind control}) that already has its own home in the structural-concepts and concept-backbone reference vocabularies (\ref{sec:concept_backbone}, \ref{sec:structural_concepts}); pooling both without discrimination would blur that distinction rather than add coverage. Three checks against real data, applied before any candidate is embedded or pooled, address this:

\textbf{A casing filter.} \texttt{domain\_terms} are recorded verbatim, in the source text's own casing. Checked directly: terms that are \emph{not} entirely lowercase are almost always genuine proper nouns (\enquote{ISKCON}, \enquote{Branch Davidians}, \enquote{UNADFI}); all-lowercase terms are almost always generic vocabulary (\enquote{secularization}, \enquote{deprogramming}, \enquote{mind control}). This filter is applied to entity anchors as well, which had never had any type filtering at all — the reason the first-generation ranking's top entity by mention count after \enquote{Scientology} was \enquote{charismatic leader} (\ref{sec:emergent_entities}), a role description rather than a named entity. One defect in the filter itself was caught and fixed before use: Python's string-is-lowercase test is false for a string with no cased characters at all, so bare years and page numbers (\enquote{2004}, \enquote{11}) were passing as though capitalized; the filter now additionally requires at least one alphabetic character.

\textbf{A cited-author stop-list.} Even after the casing filter, the single largest remaining source of noise in the literature corpus was bare surnames of the corpus's own cited authors — an in-text citation such as \enquote{(Richardson 1985)} leaves a capitalized surname indistinguishable, by casing or frequency alone, from a genuine named entity. Checked directly: \enquote{Richardson} alone had 47 mentions after the casing filter, ahead of most genuine named cult groups; \enquote{Barker} 34; \enquote{Anthony} 28. A mention-count threshold cannot separate self-citation from a genuinely recurring entity, since both present identically — a capitalized word mentioned often. The fix reads every author surname directly from the literature corpus's own bibliographic metadata (\ref{sec:corpus_construction_protocol}; 90 unique surnames across the 57 approved works) and excludes them from the candidate pool — derived from the corpus's own author list rather than a hand-typed name list, so it stays correct as the corpus grows. This is a targeted fix for self-citation within this corpus's own authorship, not a general citation detector: externally-cited classical theorists who are not themselves corpus authors (checked directly: \enquote{Weber}, \enquote{Freud}, \enquote{Wallis} all still pass) remain a documented, unaddressed residual source of noise.

\textbf{Per-corpus, not global, mention thresholds.} The original three-mentions-across-the-whole-corpus threshold (\ref{sec:emergent_entities}) is replaced with a threshold applied separately per corpus, since literature's citation-heavy volume makes a given mention count far less indicative of a genuine recurring entity than the same count in MIVILUDES's or the interview corpus's much smaller, curated text; an entity qualifies for inclusion if it clears \emph{any one} corpus's own threshold (literature $\geq$10, MIVILUDES $\geq$3, interviews $\geq$1 in the current run), with its full per-corpus mention distribution retained regardless of which corpus's threshold it cleared on.

Together these three checks recover \textbf{753} emergent-entity points from the revised archive (against the first, unrevised attempt's 21) — 603 mentioned only in literature, 55 only in interviews, 26 only in MIVILUDES, 66 shared between exactly two corpora, and 3 shared across all three (\enquote{new age}, the French \enquote{témoins de jéhovah} for Jehovah's Witnesses, and the French \enquote{ordre du temple solaire} for the Order of the Solar Temple). The top-ranked entities by mention count are, throughout, genuine named groups and movements (Scientology, Heaven's Gate, the Unification Church, Jehovah's Witnesses, the Branch Davidians, Jonestown, ISKCON, Aum Shinrikyo), in contrast to the first-generation ranking's mixture of named entities and generic nouns (\ref{sec:emergent_entities}). This is reported as a substantially improved, not a fully hand-verified, candidate list: spot-checking the top-ranked entities and the shared-anchor retrieval found no remaining obvious false positives, but the underlying filters are not a general citation detector, and the residual gap for externally-cited theorists noted above is real. A separate, small, individually-documented convention handles the rare case where a specific extracted item needs to be excluded from the shared space entirely rather than filtered by a general rule: excluded keys are recorded explicitly, by their own document/chunk identifier, and applied at pooling time — never by editing a source archive, so the archive itself remains a complete, unmodified record of what the extraction pipeline actually produced. One such exclusion is currently in effect, documented in full in \ref{sec:known_limitations}. The full revised ranking, top 100 with a per-corpus mention breakdown, is reproduced in \ref{ch:emergent_entities_v2}, alongside — not replacing — the first-generation ranking in \ref{ch:emergent_entities}.

\subsection{A Shared Cross-Corpus Space}\label{sec:shared_space}
The concept backbone's stated purpose above — a common reference frame for comparing spaces built from different corpora — is not achieved merely by embedding it alongside everything else. Each per-corpus vectorisation (\ref{sec:vectorising_scholarly_work}) reduces its own archive's dimensionality independently, projecting literature's expressions into their own space, MIVILUDES's into another, and the interview corpus's into a third; these are three unrelated coordinate systems, and a point's position in one carries no meaning relative to a point's position in another. Comparing them geometrically first requires putting every point — from every corpus, the MIVILUDES criteria, the concept backbone, and the emergent entities (\ref{sec:emergent_entities}) alike — into the same space.

Every point pooled this way also carries a \textit{point role}, distinguishing three kinds of thing the space brings together: an \textbf{expression} point is a criterion expression extracted from a text, or the MIVILUDES's own criterion text — something a source actually said (literature, MIVILUDES, interviews, MIVILUDES criteria); a \textbf{reference} point is one of three backdrop vocabularies, not itself a claim any source makes — the corpus-independent concept backbone (\ref{sec:concept_backbone}), the corpus-derived structural concepts (\ref{sec:structural_concepts}), or the ConceptNet-generalised third vocabulary (\ref{sec:conceptnet_concepts}), kept as three separate subsets throughout rather than merged into one undifferentiated \enquote{reference} category; an \textbf{emergent} point is a named entity or concept mentioned \textit{by} the corpora themselves (\ref{sec:emergent_entities}) — corpus-derived like an expression, but a recurring reference object rather than a claim. This role is orthogonal to which specific corpus a point came from, and lets later analysis ask, for instance, how expression points are distributed around a reference or emergent point, without caring which of the three expression corpora contributed them.

This is done by pooling every embedded vector produced by this thesis into a single matrix, rather than reducing each source separately: each scholarly-literature, MIVILUDES, or interview expression contributes one point — except exact-duplicate expressions within the same document (a known, previously-deferred extraction artefact, \ref{sec:vectorising_scholarly_work}: the model occasionally emits the same expression twice) and expressions under five words (typically a short, low-content fragment already captured as an entity anchor on a fuller expression), both filtered here at pooling time rather than by editing the archive — 775 duplicates and 3{,}048 short fragments removed across the three corpora on the current run; each of the seventeen MIVILUDES criteria contributes one, its French embedding only (\ref{sec:vectorising_miviludes}); each of the 3,000 concept-backbone entries contributes one; each of the 1,500 structural concepts contributes one (\ref{sec:structural_concepts}); each of the 195 hand-reviewed ConceptNet-derived reference terms contributes one (\ref{sec:conceptnet_concepts}); each of the 3,251 emergent entities contributes one (\ref{sec:emergent_entities}) — 44{,}520 points in total across eight source datasets. Each pooled corpus-expression point additionally carries its claim mode, epistemic status, and attribution tag (\ref{sec:vectorising_scholarly_work}) unchanged, and interview points their response rank, so later analysis can facet or filter the space along any of them without re-deriving them from the archive. The pooled matrix is standardised (zero mean, unit variance per dimension) before reduction: although every vector comes from the same embedding model, the sources differ substantially in register — academic prose, government French, casual interview speech, bare word-and-gloss dictionary definitions, bare named entities — and could plausibly carry different per-dimension distributions, so standardising guards against any one source dominating the fit purely as an artefact of scale. A single principal component analysis is then fit on this pooled, standardised matrix, and every point — again, from every source — is projected through that one fitted transform. The target dimensionality is not fixed in advance: the full explained-variance curve is computed first, and the smallest number of components reaching 95\% cumulative variance is chosen from it (394 components on the current corpus; the curve is fairly gradual rather than showing a sharp knee — 61\% of variance at 100 components, 90\% at 300 — so this is a real but not dramatic compression, worth bearing in mind when interpreting distances in the resulting space). As a check on whether this construction behaves sensibly, each MIVILUDES criterion's French and English raw embeddings — the latter not itself pooled into the shared space, \ref{sec:vectorising_miviludes} — are compared directly by cosine similarity before reduction: mean 0.87, median 0.90, with the lower tail (10th percentile 0.81, minimum 0.50) expected from short official phrases embedding less stably cross-lingually than length alone would suggest, rather than indicating mistranslation (verified by hand for the single pair below the 0.70 inspection threshold). As a check on the three reference vocabularies' design (\ref{sec:structural_concepts}, \ref{sec:conceptnet_concepts}), each one's centroid distance to the combined expression-corpus centroid was measured (\texttt{thesis\_corpus.analyze\_global\_structure}): 16.70 (concept backbone), 16.18 (structural concepts), 14.92 (ConceptNet-derived), against 13.20 for emergent entities — the intended ordering holds: the two corpus-anchored reference vocabularies read closer to the expression cloud than the topic-neutral concept backbone, with the hand-pruned ConceptNet set closest of the three. (These figures move slightly with every pooling-time change, since standardization is fit jointly across the whole pooled matrix before PCA — adding a new source dataset, or swapping an existing one's language, shifts every dimension's mean/variance a little, not just the changed points' own coordinates; the relative ordering has held across every rerun regardless.) This centroid comparison is not adjusted for the three vocabularies' very different sizes (3,000/1,500/195 terms) — a separate, equal-size-controlled nearest-term comparison (every set down-sampled to the smallest before computing nearest-neighbour distance, repeated with bootstrap resampling) is reported alongside it for exactly that reason, and controls for pool-size only, not for the vocabularies' differing source, curation method, or coverage.

This shared space, not any single corpus's independently-reduced one, is what any subsequent geometric comparison across corpora in this thesis is built on.

\subsection{Rebuilding the Shared Space Under the Revised Pipeline}\label{sec:shared_space_v2}

The revised extraction pipeline (\ref{sec:extraction_v2}) is, deliberately, far more conservative than the first generation: literature drops from 39{,}236 to 5{,}741 expressions (57 documents, unchanged), MIVILUDES from 914 to 108 (2 documents, unchanged), interviews from 230 to 64 (26 transcripts, unchanged) — an 72--88\% reduction per corpus, trading quantity for verified verbatim fidelity rather than a change in source material. Pooling this revised archive follows the identical procedure described above (\ref{sec:shared_space}): standardise, fit one PCA on the pooled matrix, project every point through it. With the corrected emergent-entity layer (\ref{sec:emergent_entities_v2}) included, the revised shared space holds 11{,}377 points across the same eight source datasets, reduced to 393 dimensions at the same 95\% cumulative-variance threshold — close to, though not identical to, the first-generation space's 394, since the target dimensionality is re-derived from the revised pooled matrix's own variance curve rather than fixed in advance.

Two elements of the first-generation construction are deliberately \emph{not} rebuilt from the revised archive, on the same \enquote{port, don't redo} principle: an embedding is a function of its own text, independent of which extraction pipeline happened to be running when a human selected it. The interview initial-exemplar prototype layer (\ref{sec:vectorising_scholarly_work}) keeps its original 25 manually-verified selections and their existing embeddings — 21 of the 25 no longer appear verbatim under the revised pipeline's much stricter screening, confirmed directly — re-projecting them through the newly-fitted transform rather than re-selecting or re-verifying them. The three reference vocabularies (\ref{sec:concept_backbone}, \ref{sec:structural_concepts}, \ref{sec:conceptnet_concepts}) and the seventeen MIVILUDES criteria (\ref{sec:vectorising_miviludes}) are likewise unchanged, since none is derived from the literature, MIVILUDES, or interview archives.

Because literature remains the large majority of expression points under the revised pipeline as well — 97.09\% of the three corpora's combined expression points, against 97.44\% under the first-generation pipeline — every equal-size/equal-corpus-weighted correction described elsewhere in this chapter (\ref{sec:known_limitations}) remains necessary and is applied unchanged; only the absolute size of the resulting equal-sample draw shrinks, from 204 points per corpus under the first-generation pipeline to 64 under the revised one, since the draw size is set by whichever corpus is smallest. Every bootstrapped confidence interval computed against the revised space accordingly rests on a smaller, and correspondingly noisier, resample than its first-generation counterpart, a limitation carried forward explicitly rather than left implicit.

\textbf{Individually-documented manual exclusions.} One extracted item — interview \texttt{b3-aug23-1213}, the participant's response to the opening prompt (\enquote{Quand je te dis le mot secte à quoi tu penses?}, \enquote{when I say the word `secte' [cult], what do you think of?}) — was manually excluded from the pooled space after a targeted fidelity check (below) found the transcribed response, \enquote{Sept?} (\enquote{seven?}), to almost certainly be a mishearing of \textit{secte} as \textit{sept} by the 97-year-old participant, rather than a substantive answer. It had passed every automated check in the revised pipeline (\ref{sec:extraction_v2}) — the second-model verification step correctly scored it faithful, self-contained, cult-relevant, and atomic, since it genuinely is a verbatim transcript quote — precisely because none of those checks are designed to detect that a faithfully-transcribed string carries no propositional content about its ostensible subject. The exclusion is applied only at pooling time, by an explicit, individually-documented key (the item's own document and chunk identifier), never by editing the extraction archive itself, which remains a complete and unmodified record of what the pipeline produced, this item included. This is the general mechanism, not a one-off special case: any future item requiring the same treatment is added to the same explicit exclusion list, each entry documented with the specific reason for its exclusion, rather than by loosening or special-casing a general extraction or screening rule.

\textbf{A targeted fidelity check of the revised archive.} Before relying on the revised pipeline's output, a fresh, reproducible random sample was drawn directly from each corpus's final, judge-accepted archive (100 literature, 40 MIVILUDES, 30 interview expressions; seed fixed for reproducibility) and checked two ways: automated verification that each sampled expression's text is a verbatim substring of its own recorded surrounding context, and an independent read of a subset of rows. Zero of 170 sampled rows failed the verbatim check; the independent read found expressions correctly attributed and consistently tagged throughout, with the \enquote{Sept?} case above as the one genuine exception found. This check also directly validated a specific design choice in the revised extraction screen (\ref{sec:extraction_v2}): unlike the first-generation pipeline's plain word-count floor, the revised screen does not reject a short span merely for being short, so as not to discard meaningful short lay associations (e.g.\ \enquote{Tomato cult!}, \enquote{Illuminati}). Checked directly across all three corpora: of every extracted item two words or fewer (1{,}132 in literature, 7 in MIVILUDES, 17 in interviews), all but one are genuine named entities or associations (\enquote{Scientology}, \enquote{the Moonies}, \enquote{les sectes}, \enquote{Gourou}) — the sole exception being the \enquote{Sept?} case documented above, which the screen's own faithfulness criteria could not have been expected to catch, precisely because it is faithful.

\subsection{Known Limitations}\label{sec:known_limitations}

\textbf{Corpus imbalance.} Literature's 35{,}621 pooled expression points (after filtering, \ref{sec:shared_space}) are roughly 97\% of all expression points, against MIVILUDES's 732 and interviews' 204. Any unweighted statistic computed over the pooled space — a grand centroid, "the nearest concept to the corpus as a whole" — is therefore a statistic dominated by literature's own vocabulary and register, not a balanced comparison across the three epistemologies. This is measured, not assumed: the equal-corpus-weighted grand centroid (each of the three expression corpora contributing equal total weight, regardless of point count) sits 3.98 shared-space units from the plain unweighted one (\texttt{thesis\_corpus.balanced\_analysis}, \texttt{weighted\_centroid()}). All 3-D visualisations and the PCA fit itself keep every literature point, so as not to lose that corpus's own internal structure; but any later quantitative claim about the corpus \textit{as a whole} should use either the equal-corpus-weighted centroid or the accompanying stratified-by-document literature subsample (\texttt{literature\_balanced\_sample.jsonl}, 2{,}500 points, proportional to each document's share of the corpus) rather than the raw pooled space directly.

\textbf{MIVILUDES's thin evidentiary base.} The MIVILUDES corpus is drawn from two documents — the 2022–2024 activity report and the \enquote{Comment identifier une dérive sectaire?} reference page (\ref{sec:vectorising_miviludes}). This is not, and should not be read as, a representative sample of French state framing of \enquote{cult}-adjacent phenomena in general; it is one influential operational framework, worth comparing precisely because MIVILUDES is the concrete institutional apparatus this thesis's legal/administrative epistemology centres on, not because two documents can stand in for the state's position more broadly.

The fitted scaling and PCA transformation are persisted (not merely fit and discarded) precisely so that a vector computed after this pooling step can still be placed into the same coordinate system without refitting anything — used both for the interview initial-exemplar prototype layer above (\ref{sec:vectorising_scholarly_work}) and, separately, to check whether the seventeen MIVILUDES criteria's language of embedding affects their measured position: since the criteria are pooled in French but the corpora are predominantly English, their own stored English translations were projected through this same persisted transform into the identical 394-dimensional space, rather than compared only in the raw, un-standardised embedding space or left unchecked. This lets the French-primary and English-projected representations be compared on identical footing — same coordinates, same distance metric, same corpus vectors — which a comparison against the raw embedding space alone could not offer, since that space is neither standardised nor dimensionality-reduced the way the shared space is; a supplementary raw-embedding cosine comparison is still reported alongside it as a lower-level diagnostic.

Two-dimensional projections are used two ways in this work, for different purposes. A single UMAP \autocite{mcinnesUMAPUniformManifold2020} projection and a linear PCA slice (the leading two dimensions of the same reduction described above) of the entire pooled space give a global, exploratory overview, but at over 44,000 points are too dense to read any one specific comparison off of directly. For readable figures of a specific comparison — one expression corpus against the three reference vocabularies, one MIVILUDES criterion against its nearest neighbours, the interview prototype layer against the criteria — a separate, small, curated point-set is projected on its own (again both by a fresh UMAP fit and by the same fixed linear PCA basis, so every such figure remains comparable to the others under the linear projection specifically). These smaller, focused projections are a visualisation aid for reading a specific comparison; the quantitative claims in this thesis rest on the full 394-dimensional distance/similarity computations described throughout this section, not on any two-dimensional projection of them.

\textbf{Language asymmetry.} MIVILUDES's own text is almost entirely French, while both reference point-sets it is compared against — the concept backbone (\ref{sec:concept_backbone}) and structural concepts (\ref{sec:structural_concepts}) — are English-only. This was not a hypothetical concern: the seventeen official criteria's own French/English translation pairs, embedded separately as a sanity check, show cosine similarities ranging from 0.50 to 0.95 for identical content (\ref{sec:shared_space}) — direct evidence that cross-lingual distance in this space carries real noise on top of whatever semantic signal is being measured. To close this gap, all 914 MIVILUDES expressions (before the duplicate/short-fragment filter, \ref{sec:shared_space}) were machine-translated to English (\texttt{qwen3:4b}) and re-embedded; the English translation is what is pooled into the shared space for each of the 732 surviving MIVILUDES expression points, with the original French retained as a \texttt{label\_fr} display field, mirroring the seventeen criteria's own construction (\ref{sec:vectorising_miviludes}) but inverted. Translation fidelity was checked the same way as the criteria: raw French/English embedding cosine similarity across all 914 pairs gives mean 0.90, median 0.91, with the lower tail (10th percentile 0.82, minimum 0.55) again consistent with short phrases embedding less stably cross-lingually rather than mistranslation; 12 of 914 pairs (1.3\%) fell below the 0.70 hand-inspection threshold, all from the same two source documents and confirmed on inspection to be accurate translations of genuinely short or fragmentary French phrases. Separately, comparing each English translation's length to its French source surfaced 5 of 914 (0.55\%) cases where the model returned an explanation of the phrase (e.g. \enquote{The phrase "situation de dépendance" is a French noun phrase that translates directly into English as...}) rather than a bare translation; 3 of these 5 survived into the final 732 pooled points and are left as-is rather than hand-corrected, a small, quantified, documented residual noise source consistent with this section's general policy of disclosing known extraction/translation artefacts rather than silently patching them. Interview transcripts are deliberately left in their original language regardless (\ref{sec:vectorising_scholarly_work}), since only 19\% of interviews are French, a much smaller asymmetry than MIVILUDES's near-total French text.

\textbf{Interview sample.} The 26-interview corpus (\ref{ch:survey_methodology}) is a convenience sample recruited through the researcher's own network, not a representative sample of lay usage of \enquote{cult}. This shows up directly in the material (\ref{sec:vectorising_prototypes}): one interview's \enquote{first association} response was excluded from the database for having been biased by a prior conversation with the researcher, and a separate respondent named the researcher's own academic programme a \enquote{cult} — both signs of a small, socially proximate sampling pool rather than an independent cross-section. The interview corpus is accordingly treated here as exploratory prototype data — a source of candidate structures and associations worth testing against the other two epistemologies — rather than as evidence about how \enquote{the lay public} in general uses the term.

\textbf{Two corrected metadata records.} Hand-inspection of the interview archive's \texttt{attribution} field (the mechanism separating a participant's own words from the interviewer's, \ref{sec:vectorising_scholarly_work}) surfaced one further mislabel beyond the ones already documented there: document \texttt{b1-aug12-1231}, chunk 3 (\enquote{how it's perceived as a whole}) — the interviewer's own paraphrased question — had been tagged \texttt{attribution=participant} instead of \texttt{unspecified}. Separately, manual review of the geometric-analysis toolkit's qualitative shared-anchor retrieval output surfaced a second mislabel, this time in \texttt{epistemic\_status}: document \texttt{b2-aug13-1840}, chunk 2 (\enquote{Scientology is also a cult}) had been tagged \texttt{negated}, but its full source context (\enquote{Or maybe Scientology is also a cult, maybe? Not sure about it.}) neither affirms nor denies the claim — it expresses uncertainty about it, so \texttt{speculative} is the correct tag under this project's epistemic-status scheme. Both corrected by hand, consistent with this project's existing practice for individually-identified mislabels of this kind. The shared space and interview-prototype layer were rebuilt after each correction; verified both times to produce byte-identical \texttt{shared\_space\_vector}s for all 44{,}520 points, since each fix touches only a single metadata field, never an embedding. Because both \texttt{negated} and \texttt{speculative} are excluded by the \texttt{asserted\_qualified} epistemic-status filter (below), the second correction changed no filtered-pool membership and left every filtered numeric output byte-identical to before it. The geometric-analysis toolkit's current authoritative outputs are run \texttt{20260908-175344} (baseline, all 9 modules, unfiltered) and run \texttt{20260908-180406} (\texttt{analyze\_global\_structure}/\texttt{analyze\_cluster\_structure} only, restricted to \texttt{asserted}/\texttt{qualified} statements, as an epistemic-status robustness check) — both supersede the intermediate \texttt{20260908-170432}/\texttt{20260908-171742} runs (attribution-corrected but pre-dating the epistemic-status correction) and the earlier \texttt{20260907-002832} run, all of which (along with \texttt{review\_pass1}/\texttt{smoketest5}/\texttt{test\_ie\_fast}) are kept on disk as superseded historical record rather than deleted. These corrections, and the archive and analysis runs listed here, belong to the first-generation pipeline's frozen developmental baseline (\ref{sec:vectorising_scholarly_work}); the corpus produced by the revised extraction pipeline (\ref{sec:extraction_v2}) was built independently from the same underlying source texts and is not affected by them.

\subsection{Vectorising Prototypes}\label{sec:vectorising_prototypes}

- I realised that the data I had gathered is a mixed bag of answers of different length and depth, mixing multiple examplars "named entities" ('Manson Family), definition attempts (document), free associations (one that came out a lot is the link between cult and religion, which can be either an association - instagram interview 7 - and a named entity in INTERVIEW — Aug 9, 20:48 ).
- This is a problem because this raises the question of which object I will use in my gemotrical analysis. I have several options:
Methods that still keep the answers without post-processing:
- Concatenate the raw answers of the interviewees and use them as is.
	- 
	- multi-entities input, for example, multiple cults objects can be mentioned, and multiple dimensions are often given. can result in quite noisy output.
	- treats cults and dimensions together
	- doesn't introduce any biaises a posteriori
	- does keep the biases that the conversational aspect of the survey brings
- Use interviewer questions as seperators, and vectorise each part seperately
	- would have the advantages and drawbacks as the previous methods, but give a more granular output

Methods that would post-process in order to get a cleaner output:

- Extract objects from transcripts. 
- Objects that I can extract: (or with an LLM) all the 'named entities', 'abstract groups' (group of several cults without a specific name, higher-order category), 'cult-dimensions', and finally the 'first association' deemed as prototypical.
- I have the 

- Introduces a biaise because it eliminates the context and the words around, and also eliminates how a person perceived the group. For example people can mention "the Moonies" without knowing what it is, but remeber they heard about it (INTERVIEW — Aug 13, 18:32).
---
Which method did I choose:
- Extraction Methodology:
	- Used an LLM: Claude 5 and supervised for, no introduction of biaises with the instruction of no rephrasing or introduction of additional terms.
	- Extracted the following objects in order of appearance:
		- 'Named Entities of cults'
		- 'Cult dimensions'
		- 'First association'
	- This poses problems: "Cult of tomatoes" or "design and computation", are quite personal entities that are not know. This was comfirmed by probing the model associate with the used embeddings model. 
	
	I probed the model with every single (see appendix). For the case that these groups were unknown, I decided to attach a short context with what this group, as defined by the respondant. 
	 


Make a simplified diagram of the methodology of extraction. 


Notes: 
- I have decided to exclude from my database the "first association" of "INTERVIEW — Aug 6, 19:55" because the answer has been biaised by a talk we previously had and awereness on my research.
- There was also an answer to the survey that desingated as a 'cult' the very academic program in which this research project is being written. 
One survey response explicitly labeled the very academic program hosting this research project as a “cult.” In the interest of intellectual honesty, I retained this response under the categories “first association” and “named entities and cult dimensions,” but I removed the associated demographic data and transcript to minimize any risk of identification or harm. This achieved with the respondant's consent.

- See ethics section

% Qualitative survey data is difficult to extract dimensions from manually --
% problems of arbitrary interpretation (e.g. category assignment) and of
% matrix growth (every dimension must be accounted for, for every object).

\section{Embeddings}
% Using pre-trained embeddings on a larger corpus instead of extracting
% dimensions / building a space from scratch: query each object's vector
% coordinate in a large latent space. Embeddings provide one reproducible
% operationalisation of semantic proximity, while introducing their own
% theoretical and methodological assumptions.

\section{Representation}
Might leave this for the practical project part.
% Comparing representation/dimensionality-reduction techniques (e.g. t-SNE,
% UMAP) and which to use; cosine similarity or other distance measures.

\section{Analysis}\label{sec:analysis_methods}

Once every corpus, reference vocabulary, and emergent-entity set is pooled into the shared space (\ref{sec:shared_space}, \ref{sec:shared_space_v2}), every geometric technique used elsewhere in this thesis is one of a small number of reusable operations, described here once rather than re-derived at each point of use. All operate on the full reduced-dimensional space (393/394 components, \ref{sec:shared_space}) unless explicitly noted as a two-dimensional projection, per the standing rule that two-dimensional pictures are a reading aid and never themselves the evidence for a quantitative claim (\ref{sec:known_limitations}).

\subsection{Centroids, dispersion, and typicality}
A group's \textit{centroid} is the mean of its member vectors — the group's own \enquote{centre of gravity} in the shared space — and its \textit{dispersion} the mean Euclidean distance from each member to that centroid. Comparing two groups' centroids (e.g.\ two corpora's own average position) reduces a many-point comparison to a single distance and a single cosine similarity, reported side by side throughout rather than collapsed into one number, since the two measures do not always agree. Because the three expression corpora differ enormously in size (\ref{sec:known_limitations}), any statistic computed over the pooled space unweighted is dominated by literature's own vocabulary and register; an \textit{equal-corpus-weighted} variant (each corpus contributing equal total weight regardless of point count) and an \textit{equal-size} or \textit{equal-$n$} variant (every group down-sampled to the smallest before comparing, repeated under bootstrap resampling for a confidence interval) are used throughout wherever a comparison could otherwise be an artefact of corpus size rather than of the underlying geometry — never silently substituted for the unweighted figure, but reported alongside it.

A point's \textit{typicality} with respect to a given group is operationalised purely as its distance to that group's own centroid: closer is \enquote{more typical} of that group's average position, farther is \enquote{more atypical} — a geometric fact about this embedding space under this model, not a claim about cognitive prototypicality or real-world category membership. A \textit{borderline} point between two groups is one whose distances to each group's centroid are nearly equal; this identifies a point's geometric position as ambiguous between two centroids, not that the underlying claim it expresses is itself contested or unclear.

\subsection{Nearest-neighbour retrieval and cluster composition}
Given an arbitrary query vector (a centroid, a single point, a reference term), every candidate in a chosen pool can be ranked by ascending Euclidean distance, with cosine similarity reported as a secondary descriptive field never used to select or order results. This single ranking operation underlies both a \enquote{$k$ nearest neighbours} retrieval (the closest $k$ candidates) and, by simply reversing the same ranking, a \enquote{$k$ farthest} one — always the same underlying sort, never two independently-computed rankings that could drift out of sync.

Applied to whole groups rather than a single query, the same operation yields \textit{$k$-nearest-neighbour composition}: for a sample of query points from one source, what fraction of their own $k$ nearest neighbours come from each of the sources present in the pool, averaged and bootstrap-resampled for a confidence interval. A \textit{silhouette score} summarises the same underlying question — whether a proposed grouping (e.g.\ the three expression corpora) corresponds to genuinely separate geometric clusters — as a single number, conventionally between $-1$ and $1$, with values near zero indicating the proposed groups do not correspond to distinct clusters and values near one indicating they do.

\subsection{Entity topology: density-based clustering and centrality}\label{sec:entity_topology_methods}
The emergent-entity layer (\ref{sec:emergent_entities}, \ref{sec:emergent_entities_v2}) is additionally analysed for its own internal spatial structure, independent of any comparison to the expression corpora. Density-based clustering (HDBSCAN \autocite{mcinnesHdbscanHierarchicalDensity2017}, via \texttt{scikit-learn} \autocite{JMLR:v12:pedregosa11a}, chosen over a fixed-$k$ method such as $k$-means specifically because it does not require the number of clusters to be specified in advance, and explicitly leaves points that do not sit inside a sufficiently dense neighbourhood as unclustered \enquote{noise} rather than force-assigning every point to some cluster) is run directly on the entity vectors in the full reduced-dimensional space, at a minimum-cluster-size threshold chosen to keep the smallest meaningful groupings (three members) without treating every incidental pair of nearby points as its own cluster. Each individual entity's own local density is separately reported as the Euclidean distance to its $k$-th nearest other entity ($k=5$): a smaller distance indicates a denser immediate neighbourhood, independent of which (if any) cluster the entity was assigned to. Each cluster's own centrality is the distance from that cluster's member centroid to the shared space's equal-corpus-weighted grand centroid (above) — the same distance-to-centroid framing used for typicality, applied here to a cluster rather than a single point, so that \enquote{central} and \enquote{peripheral} clusters are identified on a consistent basis throughout this thesis. As with typicality generally, cluster membership, local density, and centrality are purely geometric facts about this space under this model, not a claim that clustered entities are conceptually similar in some deeper sense, nor that an unclustered (\enquote{noise}) entity is unimportant.

\subsection{Two-dimensional partitioning: Voronoi tessellation and density gaps}\label{sec:voronoi_methods}
To read the entity layer's spatial structure as a picture rather than a table, the entity vectors are projected to two dimensions by UMAP \autocite{mcinnesUMAPUniformManifold2020}, fit once on the entity population alone; a chosen set of \textit{seed} points (e.g.\ the seventeen MIVILUDES criteria, the three corpus centroids, or the entity clusters' own centroids) is then projected into that same fitted two-dimensional layout and used to build a Voronoi tessellation — a partition of the plane into regions, one per seed, each region containing exactly the points closer to that seed than to any other. Several seed strategies can be tried against the identical two-dimensional entity layout and compared directly, since the layout itself does not change between them (verified, for the seed sets used in this thesis, to produce numerically identical entity coordinates across separate fitting runs, UMAP being deterministic given a fixed random seed) — only which seed set the tessellation is built from differs.

This technique carries one documented limitation, load-bearing enough to record here rather than only alongside whichever results it affects: a seed drawn from the same population UMAP was fit on projects sensibly into the resulting two-dimensional layout, but a seed from an unrelated population (e.g.\ the MIVILUDES criteria or the corpus centroids, neither of which are entity vectors) only enters the fitted layout afterward, via UMAP's out-of-sample transform — a well-documented behaviour of the underlying algorithm in which points not part of the original fit can be compressed into an unrepresentative region of the projection, regardless of their actual position in the full-dimensional space. Whether a given seed strategy's two-dimensional tessellation is legible is accordingly not itself evidence about the seeds' true geometric relationship to the entities; the underlying full-dimensional nearest-neighbour retrieval (above) is unaffected by this projection-specific limitation and is what any quantitative claim in this thesis relies on.

The same two-dimensional entity layout is also used to identify candidate coverage gaps: binned into a fixed grid over its own bounding box, a cell is a candidate gap if it contains no entities \textit{and} sits inside the convex hull of the entities' own two-dimensional positions — the second condition specifically excluding a cell that is merely empty space at the edge of the projection, so that a candidate gap is a region genuinely surrounded by occupied space rather than simply outside the data's own extent.

\subsection{Publication-date effects}\label{sec:publication_date_methods}
For the scholarly-literature corpus specifically (\ref{sec:corpus_construction_protocol}), each expression's document is additionally associated with a publication year, recovered directly from the document identifier itself (\ref{sec:pdf_extraction}) via pattern matching for a four-digit year token, rather than through the corpus's separate bibliographic metadata record, whose own identifier scheme does not share a join key with the extraction pipeline's document identifiers. This is applied to the literature corpus only: neither MIVILUDES (two documents, no structured per-document date) nor the interview corpus (dated within one narrow collection window, not a multi-year span) supports a meaningful publication-date analysis on the same basis. Two measures are computed from the resulting year-per-point association: a rank-based (Spearman) correlation between publication year and a point's distance to the literature corpus's own centroid, chosen over a linear correlation coefficient since it assumes only a monotonic, not necessarily linear, relationship; and a period-level comparison, splitting the corpus into year terciles using the corpus's own year distribution (its 33rd and 66th percentiles) rather than a fixed external calendar grid, so that the resulting periods reflect this corpus's actual publication-year spread rather than an arbitrary division of the calendar. Each period's own centroid can then be compared to the others', and to the shared space's other reference points (e.g.\ the MIVILUDES criteria), by the same nearest-neighbour and pairwise-distance operations described above.

\subsection{Supplementary technical reports}
Every technique described in this section is implemented as a standalone, independently-runnable module in a reusable analysis toolkit (\texttt{thesis\_corpus/analyze\_*.py}, \texttt{generate\_*.py}), rather than as one-off analysis code, so that any figure or table in this thesis can be regenerated, and its exact computation inspected, from the same source. Two self-contained technical reports, deliberately kept separate from this chapter and from the Results chapter, walk through every module's output in full, plain-language detail, with every raw table and figure the toolkit produces: \texttt{04\_Appendix/Geometric\_Analysis\_Draft/} for the first-generation pipeline's space, and \texttt{04\_Appendix/Geometric\_Analysis\_Draft\_v2/} for the revised one (\ref{sec:shared_space_v2}). Both are explicitly working documents — neutral, descriptive material for writing the Results chapter from, not the Results chapter itself — and are referenced here as the primary supplementary record of this chapter's methods actually being applied, rather than reproduced inline.
